\section{Implementazione di Insiemi e Multi-Insiemi (Lez. 09)}

\subsection{Usare gli Oggetti come Insiemi (Sets)}
In JavaScript, possiamo usare gli Oggetti per simulare gli insiemi matematici.
L'idea è semplice: usiamo le chiavi dell'oggetto per rappresentare gli elementi. Il valore non importa (spesso mettiamo 1 o true).
Questo garantisce che \textbf{non ci siano duplicati} (una chiave in un oggetto è unica).

\begin{concept}{Unicità delle Chiavi}
L'idea è semplice: usiamo le chiavi dell'oggetto per rappresentare gli elementi. Il valore non importa (spesso mettiamo 1 o true).
Questo garantisce che \textbf{non ci siano duplicati} (una chiave in un oggetto è unica).
\end{concept}
\begin{jsbox}
let frutta = {
    "mela": 1,
    "pera": 1,
    "banana": 1
};

// Verifica Appartenenza (in)
if ("mela" in frutta) {
    console.log("C'e' la mela!");
}

// Aggiungere (basta assegnare)
frutta["kiwi"] = 1;

// Rimuovere (delete)
delete frutta["pera"];
\end{jsbox}

\subsection{Operazioni sugli Insiemi: Unione e Intersezione}
Possiamo scrivere semplici funzioni per combinare due insiemi.

\subsubsection{Unione (+)}
\begin{concept}{Somma Logica (OR)}
Prendiamo due oggetti e ne creiamo un terzo che ha le chiavi di entrambi.
\end{concept}
\begin{jsbox}
function unione(A, B) {
    let C = {};
    // Copio tutte le chiavi di A
    for (let k in A) C[k] = 1;
    // Copio tutte le chiavi di B (se ci sono già, sovrascrivo, non fa nulla)
    for (let k in B) C[k] = 1;
    return C;
}
\end{jsbox}

\subsubsection{Intersezione (Solo i comuni)}
\begin{concept}{Prodotto Logico (AND)}
Tengo solo gli elementi che esistono in entrambi gli insiemi.
\end{concept}
\begin{jsbox}
function intersezione(A, B) {
    let C = {};
    for (let k in A) {
        // Se c'è anche in B, lo tengo
        if (k in B) {
            C[k] = 1;
        }
    }
    return C;
}
\end{jsbox}

\subsection{Multi-Insiemi (Contatori)}
A volte vogliamo contare \textit{quante volte} appare un elemento (es. carrello della spesa: 3 mele).
In questo caso, il valore associato alla chiave non sarà più 1, ma il numero di occorrenze.

\begin{jsbox}
let carrello = {
    "mela": 3,  // 3 mele
    "latte": 2, // 2 bottiglie
    "pane": 1
};
\end{jsbox}

\subsubsection{Da Array a Multi-Insieme}
\begin{concept}{Istogramma di Frequenza}
Un esercizio classico: dato un array di parole, contare quante volte appare ciascuna.
\end{concept}
\begin{jsbox}
function contaParole(lista) {
    let contatore = {};
    for (let parola of lista) {
        if (parola in contatore) {
            contatore[parola] += 1; // Incremento
        } else {
            contatore[parola] = 1; // Inizializzo
        }
    }
    return contatore;
}

let input = ["mela", "pera", "mela"];
console.log(contaParole(input)); // { mela: 2, pera: 1 }
\end{jsbox}
